% ============================================================
%  LUAN VAN THAC SI - KHUNG (SKELETON)
%  Bien dich bang XeLaTeX (khong dung pdfLaTeX vi dung fontspec)
%    xelatex luan_van_khung.tex
%    xelatex luan_van_khung.tex   (chay 2 lan de muc luc dung)
% ============================================================

\documentclass[a4paper,12pt]{report}

% ---------- Font & ngon ngu (Unicode tieng Viet qua XeLaTeX) ----------
\usepackage{fontspec}
\setmainfont{Times New Roman}

% ---------- Trang & khoang cach (chuan trinh bay luan van) ----------
\usepackage[a4paper,top=3cm,bottom=3cm,left=3.5cm,right=2cm]{geometry}
\usepackage{setspace}
\onehalfspacing

% ---------- Cong cu toan hoc, bang, hinh ----------
\usepackage{amsmath,amssymb}
\usepackage{graphicx}
\usepackage{booktabs}       % ke bang dep (\toprule \midrule \bottomrule)
\usepackage{longtable}      % bang dai qua 1 trang
\usepackage{caption}
\usepackage{chngcntr}       % cho phep \counterwithin

% Danh so Bang/Hinh theo chuong: Bang 4.1, Hinh 3.1, ...
\counterwithin{table}{chapter}
\counterwithin{figure}{chapter}

% ---------- Tieu de chuong dang "CHUONG 1. TEN CHUONG" ----------
\usepackage{titlesec}
\titleformat{\chapter}[hang]
  {\normalfont\bfseries\Large}{CH\MakeUppercase{ương} \thechapter.}{0.8em}{\MakeUppercase}
\titlespacing*{\chapter}{0pt}{0pt}{2em}

% ---------- Doi ten cac muc tu dong sang tieng Viet ----------
\renewcommand{\contentsname}{MỤC LỤC}
\renewcommand{\listtablename}{DANH MỤC BẢNG}
\renewcommand{\listfigurename}{DANH MỤC HÌNH}
\renewcommand{\bibname}{TÀI LIỆU THAM KHẢO}

% ---------- Sieu lien ket (dat cuoi cung trong phan khai bao) ----------
\usepackage[hidelinks]{hyperref}
\hypersetup{
  pdftitle={Hieu Lan duong va Bien bao Giao thong theo Ngu nghia su dung Mo hinh Ngon ngu Lon},
  pdfauthor={Do Minh Hieu},
  bookmarksnumbered=true,
}

% ============================================================
\begin{document}

% ------------------------------------------------------------
% TRANG BIA
% ------------------------------------------------------------
\begin{titlepage}
\centering
\Large \textbf{VIỆN QUẢN TRỊ VÀ CÔNG NGHỆ FPT}\\[2cm]

\Large \textbf{HIỂU LÀN ĐƯỜNG VÀ BIỂN BÁO GIAO THÔNG THEO NGỮ NGHĨA\\
SỬ DỤNG MÔ HÌNH NGÔN NGỮ LỚN\\
ĐỂ HỖ TRỢ RA QUYẾT ĐỊNH LÁI XE}\\[0.5cm]

\normalsize \textit{(Semantic Lane and Traffic Sign Understanding Using\\
Large Language Models for Driving Decision Support)}\\[2.5cm]

\Large \textbf{LUẬN VĂN THẠC SĨ}\\[2cm]

\normalsize
\begin{tabular}{ll}
Chuyên ngành: & Kỹ thuật Phần mềm \\[0.3cm]
Học viên thực hiện: & Đỗ Minh Hiếu \\[0.3cm]
Người hướng dẫn khoa học: & Tiến sĩ Đoàn Nhật Quang \\
\end{tabular}

\vfill
Hà Nội, 2026
\end{titlepage}

% ------------------------------------------------------------
% FRONT MATTER (danh so trang La Ma)
% ------------------------------------------------------------
\pagenumbering{roman}

\chapter*{LỜI CAM ĐOAN}
\addcontentsline{toc}{chapter}{LỜI CAM ĐOAN}
% TODO: dan noi dung Loi cam doan tu ban thao .md vao day.

\chapter*{LỜI CẢM ƠN}
\addcontentsline{toc}{chapter}{LỜI CẢM ƠN}
% TODO: dan noi dung Loi cam on tu ban thao .md vao day.

\tableofcontents

\listoftables
\addcontentsline{toc}{chapter}{DANH MỤC BẢNG}

\listoffigures
\addcontentsline{toc}{chapter}{DANH MỤC HÌNH}

\chapter*{DANH MỤC TỪ VIẾT TẮT}
\addcontentsline{toc}{chapter}{DANH MỤC TỪ VIẾT TẮT}
% TODO: dan bang tu viet tat (ADAS, VLM, LLM, MLLM, UFLD-v2, SCNN, YOLO,
% TT100K, MAE, IoU, F1, API, NIM, JSON, RQ, QCVN) tu ban thao .md vao day,
% dung moi truong tabular/longtable, vi du:
%
% \begin{longtable}{p{2.5cm} p{5cm} p{7cm}}
% \toprule
% \textbf{Từ viết tắt} & \textbf{Tiếng Anh} & \textbf{Giải thích} \\
% \midrule
% \endhead
% ADAS & Advanced Driver Assistance System & Hệ thống hỗ trợ lái xe tiên tiến \\
% ...
% \bottomrule
% \end{longtable}

\chapter*{TÓM TẮT}
\addcontentsline{toc}{chapter}{TÓM TẮT}
% TODO: dan noi dung Tom tat tu ban thao .md vao day.
%
% \textbf{Từ khóa}: % TODO

% ------------------------------------------------------------
% NOI DUNG CHINH (danh so trang A-rap)
% ------------------------------------------------------------
\pagenumbering{arabic}

% ============================== CHUONG 1 ==============================
\chapter{GIỚI THIỆU}

\section{Bối cảnh và động lực}
% TODO

\section{Mục tiêu nghiên cứu}
% TODO

\section{Phạm vi dữ liệu}
% TODO

\section{Câu hỏi nghiên cứu}
% TODO

\section{Đóng góp chính}
% TODO

\section{Cấu trúc luận văn}
% TODO

% ============================== CHUONG 2 ==============================
\chapter{TỔNG QUAN VÀ CÁC CÔNG TRÌNH LIÊN QUAN}

\section{Phát hiện làn đường (Lane Detection)}
% TODO

\section{Phát hiện biển báo giao thông (Traffic Sign Detection)}
% TODO

\section{Mô hình ngôn ngữ lớn đa phương thức cho hỗ trợ quyết định lái xe}
% TODO

\section{Đánh giá chất lượng output ngôn ngữ tự nhiên bằng LLM-as-a-Judge}
% TODO

\section{Khoảng trống nghiên cứu}
% TODO

% ============================== CHUONG 3 ==============================
\chapter{PHƯƠNG PHÁP LUẬN}

\section{Kiến trúc hệ thống tổng thể}
% TODO
%
% Vi du chen Hinh 3.1 (dung goi trong package graphicx):
% \begin{figure}[htbp]
%   \centering
%   \includegraphics[width=0.9\textwidth]{hinh/pipeline.png}
%   \caption{Kiến trúc tổng thể của pipeline bốn tầng.}
%   \label{fig:pipeline}
% \end{figure}

\section{Dữ liệu}
% TODO

\section{Phân tích ngữ nghĩa làn đường}
% TODO
%
% Vi du cong thuc hien thi (display equation):
% \begin{equation}
%   y_{ref} = 0{,}95 \times \text{image\_height}
% \end{equation}
%
% Vi du he phuong trinh dieu kien (cases):
% \begin{equation}
%   p_{between} =
%   \begin{cases}
%     0{,}5 & \text{nếu xe nằm giữa hai đường biên} \\
%     1     & \text{ngược lại}
%   \end{cases}
% \end{equation}

\section{Cấu trúc JSON ngữ nghĩa gửi cho LLM}
% TODO

\section{Thiết kế prompt cho tầng suy luận}
% TODO

\section{Lựa chọn mô hình cho tầng suy luận}
% TODO

\section{Phương pháp luận đánh giá}
% TODO

% ============================== CHUONG 4 ==============================
\chapter{KẾT QUẢ VÀ BÀN LUẬN}

\section{Độ chính xác module hiểu làn đường (CULane)}
% TODO
%
% Vi du chen Bang 4.1 (dung goi trong package booktabs):
% \begin{table}[htbp]
%   \centering
%   \caption{Độ chính xác module hiểu làn đường trước và sau khi sửa lỗi off-by-one (N=200).}
%   \label{tab:bang41}
%   \begin{tabular}{lcc}
%     \toprule
%     \textbf{Chỉ số} & \textbf{Trước khi sửa lỗi} & \textbf{Sau khi sửa lỗi (N=200)} \\
%     \midrule
%     Lane count -- Accuracy & 11,1\% & \textbf{69,0\%} \\
%     \bottomrule
%   \end{tabular}
% \end{table}

\section{Kiểm chứng độc lập trên dữ liệu real-life}
% TODO

\section{Module biển báo giao thông trên CULane}
% TODO

\section{So sánh mô hình LLM cho tầng suy luận}
% TODO

\section{Kết quả chính: đóng góp của JSON ngữ nghĩa}
% TODO

\section{Kiểm chứng độ tin cậy của phương pháp đánh giá}
% TODO

\section{Hạn chế: rủi ro trùng lặp dữ liệu (data leakage)}
% TODO

\section{Kiểm chứng khả năng tự nhận diện làn đường của VLM}
% TODO

\section{Bàn luận: cơ chế đóng góp thực sự của JSON ngữ nghĩa}
% TODO

% ============================== CHUONG 5 ==============================
\chapter{KẾT LUẬN}

\section{Tóm tắt đóng góp}
% TODO

\section{Trả lời các câu hỏi nghiên cứu}
% TODO

\section{Hạn chế}
% TODO

\section{Hướng phát triển tiếp theo}
% TODO

% ------------------------------------------------------------
% TAI LIEU THAM KHAO
% ------------------------------------------------------------
\begin{thebibliography}{99}

\bibitem{qin2022ufldv2}
Z. Qin, P. Zhang, and X. Li, ``Ultra fast deep lane detection with hybrid anchor driven ordinal classification,'' \textit{IEEE Trans. Pattern Anal. Mach. Intell.}, 2022. [Online]. Available: \url{https://arxiv.org/abs/2206.07389}

\bibitem{xu2024drivegpt4}
Z. Xu, Y. Zhang, E. Xie, Z. Zhao, Y. Guo, K. K. Y. Wong, Z. Li, and H. Zhao, ``DriveGPT4: Interpretable end-to-end autonomous driving via large language model,'' \textit{IEEE Robot. Autom. Lett.}, vol. 9, no. 10, pp. 8186--8193, Oct. 2024.

\bibitem{sima2024drivelm}
C. Sima, K. Renz, K. Chitta, L. Chen, H. Zhang, C. Xie, J. Bei{\ss}wenger, P. Luo, A. Geiger, and H. Li, ``DriveLM: Driving with graph visual question answering,'' in \textit{Proc. Eur. Conf. Comput. Vis. (ECCV)}, 2024.

\bibitem{shao2024lmdrive}
H. Shao, Y. Hu, L. Wang, S. L. Waslander, Y. Liu, and H. Li, ``LMDrive: Closed-loop end-to-end driving with large language models,'' in \textit{Proc. IEEE/CVF Conf. Comput. Vis. Pattern Recognit. (CVPR)}, 2024.

\bibitem{zheng2023mtbench}
L. Zheng, W.-L. Chiang, Y. Sheng, S. Zhuang, Z. Wu, Y. Zhuang, Z. Lin, Z. Li, D. Li, E. P. Xing, H. Zhang, J. E. Gonzalez, and I. Stoica, ``Judging LLM-as-a-judge with MT-Bench and Chatbot Arena,'' in \textit{Proc. Adv. Neural Inf. Process. Syst. (NeurIPS)}, 2023.

\bibitem{pan2018culane}
X. Pan, X. Zhan, J. Shi, P. Luo, X. Wang, and X. Tang, ``Spatial as deep: Spatial CNN for traffic scene understanding,'' in \textit{Proc. AAAI Conf. Artif. Intell.}, 2018.

\bibitem{zhang2021survey}
Y. Zhang, Z. Lu, X. Zhang, J.-H. Xue, and Q. Liao, ``Deep learning in lane marking detection: A survey,'' \textit{IEEE Trans. Intell. Transp. Syst.}, 2021.

\bibitem{zakaria2023slr}
N. J. Zakaria, M. I. Shapiai, R. A. Ghani, M. N. M. Yassin, M. Z. Ibrahim, and N. Wahid, ``Lane detection in autonomous vehicles: A systematic review,'' \textit{IEEE Access}, vol. 11, pp. 3729--3765, 2023.

\bibitem{zhu2016tt100k}
Z. Zhu, D. Liang, S. Zhang, X. Huang, B. Li, and S. Hu, ``Traffic-sign detection and classification in the wild,'' in \textit{Proc. IEEE Conf. Comput. Vis. Pattern Recognit. (CVPR)}, 2016.

\bibitem{hong2019rules}
J. Hong, B. Sapp, and J. Philbin, ``Rules of the road: Predicting driving behavior with a convolutional model of semantic interactions,'' in \textit{Proc. IEEE/CVF Conf. Comput. Vis. Pattern Recognit. (CVPR)}, 2019, pp. 8454--8462.

\bibitem{shaw2025saferoute}
A. K. Shaw, C. K. Sah, X. Lian, A. S. Baig, T. Wen, K. Jiang, M. Yang, D. Yang, and L. Zhang, ``SafeRoute: Enhancing traffic scene understanding via a unified deep learning and multimodal LLM,'' in \textit{Proc. IEEE/CVF Int. Conf. Comput. Vis. Workshops (ICCVW)}, 2025.

\bibitem{sah2025advancing}
C. K. Sah, A. K. Shaw, X. Lian, A. S. Baig, T. Wen, K. Jiang, M. Yang, and D. Yang, ``Advancing autonomous vehicle intelligence: Deep learning and multimodal LLM for traffic sign recognition and robust lane detection,'' \textit{arXiv:2503.06313}, 2025.

\bibitem{kim2025dscllm}
S. Kim, J. Jin, S. Hong, D. Ka, H. Kim, and B. Noh, ``DSC-LLM: Driving scene context representation-based trajectory prediction framework with risk factor reasoning using LLMs,'' \textit{Sensors}, vol. 25, no. 23, Art. no. 7112, 2025, doi: 10.3390/s25237112.

\bibitem{ferrag2025review}
M. A. Ferrag, N. Tihanyi, and M. Debbah, ``From LLM reasoning to autonomous AI agents: A comprehensive review,'' \textit{arXiv:2504.19678}, 2025.

\bibitem{kim2026augmenting}
J. P. Kim, ``Augmenting human evaluation with LLM judges: How many human reviews do you need?'' \textit{arXiv:2605.16354}, 2026.

\bibitem{saha2026budget}
A. Saha, A. Wagde, and B. Kveton, ``LLM-as-judge on a budget,'' \textit{arXiv:2602.15481}, 2026.

\bibitem{pan2024human}
Q. Pan, Z. Ashktorab, M. Desmond, M. Santill\'an Cooper, J. Johnson, R. Nair, E. Daly, and W. Geyer, ``Human-centered design recommendations for LLM-as-a-judge,'' in \textit{Proc. 1st Hum.-Centered Large Lang. Model Workshop (HuCLLM)}, 2024.

\bibitem{zhang2025ratt}
J. Zhang, X. Wang, W. Ren, L. Jiang, D. Wang, and K. Liu, ``RATT: A thought structure for coherent and correct LLM reasoning,'' in \textit{Proc. AAAI Conf. Artif. Intell.}, vol. 39, no. 25, pp. 26733--26741, 2025.

\end{thebibliography}

\end{document}
